\clearpage
\thispagestyle{empty}
\SourceLeaf{381}\pdfbookmark[0]{Störungen der Ceres}{ceres}
\vspace*{1.6in}
\begin{center}{\Large STÖRUNGEN DER CERES.\par}\end{center}
\clearpage
\fancyhead[L]{\small STÖRUNGEN DER CERES}
\SourceLeaf{383}\pdfbookmark[1]{Briefwechsel und Nachlass}{ceres-briefwechsel}
\begin{center}
{\large BRIEFWECHSEL UND NACHLASS.\par}
\bigskip
[I. ERSTE METHODE.]\par
[OCTOBER—DECEMBER 1802.]\par
\bigskip
\pdfbookmark[2]{1. Briefe, 12. und 26. October 1802; 28. Dezember 1815}{ceres-001}[1.]\par
\medskip
\textsc{Gauss} an \textsc{Olbers}, Braunschweig, 12. October 1802.
\end{center}
\noindent\hspace*{\parindent}.\quad.\quad.\quad.\quad.\quad Ich habe angefangen, mich mit der Verbesserung der {\AstroSymbols ⚳}-Bahn zu beschäftigen. Ich habe erstlich nach den VII. Elementen die Störungen berechnet, aber bloss erst die erste Potenz der Excentricität in Betracht gezogen, weil der Bogen noch zu klein ist, um die beträchtliche Arbeit der Berechnung der höhern Potenzen zu vergüten. Ich habe einige beträchtliche Unterschiede mit \textsc{Oriani}. Mit diesen Störungen habe ich bereits eine neue Bahn, und nach dieser neuen Bahn wiederum die Störungen aufs neue bestimmt. Ich habe in allem 54 Gleichungen. Mit diesen Störungen, die von den erstern nicht \textls[180]{sehr} verschieden sind, berechne ich morgen oder übermorgen eine neue Bahn, das Resultat schicke ich Ihnen dann zusammen. Die gefundenen Störungen sind in der Länge:
\begin{center}
\begingroup\footnotesize\setlength{\tabcolsep}{4pt}\renewcommand{\arraystretch}{1.25}
\begin{tabular}{l|c|c|c|c|c}
 & 1801 Jan. 1. & Febr. 11. & Dec. 7. & 1802 März 17. & Mai 24.\\
\hline
nach d. 1. Form. & \(-428''{,}72\) & \(-358''{,}23\) & \(+646''{,}91\) & \(+954''{,}39\) & \(+1083''{,}37\)\\
»\quad »\quad 2.\quad » & \(-415{,}82\) & \(-346{,}12\) & \(+641{,}41\) & \(+945{,}31\) & \(+1069{,}42\)\\
\end{tabular}\endgroup
\end{center}
\par
\noindent\hspace*{\parindent}In der Breite ist der Unterschied ganz unbedeutend, beim Radius vector bin ich noch nicht fertig. — In der Zeit, dass ich nach den Störungsformeln den numerischen Werth berechne, könnte ich mehr wie eine Bahnbestimmung \SourceLeaf{384}machen. Übrigens hoffe ich, alle bisher vorhandenen Beobachtungen durchgehends fast haarscharf darstellen zu können, und darf so auf künftiges Jahr noch eine sehr gute Übereinstimmung erwarten. — Bei der Ceres wird man nun für die Störungen wohl noch eine brauchbare Convergenz erhalten, und ich habe schon einige Ideen zu einer bequemen Abkürzung und Zusammenziehung der Tafeln entworfen, aber bei der {\AstroSymbols ⚴} wird sicher die Theorie noch einer Art von Revolution bedürfen.\quad.\quad.\quad.\quad.\quad.
\par\bigskip
\begin{center}\textsc{Gauss} an \textsc{Olbers}, Braunschweig, 26. October 1802.\end{center}
\noindent\hspace*{\parindent}.\quad.\quad.\quad.\quad.\quad Die in meinem letzten Briefe erwähnten Rechnungen über die {\AstroSymbols ⚳} habe ich beendigt.\quad.\quad.\quad.\quad.\quad.
\par\medskip
\begin{center}\rule{.27\linewidth}{.4pt}\end{center}
\par\bigskip
\begin{center}\textsc{Gauss} an \textsc{Gerling}, Göttingen, 28. Dezember 1815.\end{center}
\noindent\hspace*{\parindent}.\quad.\quad.\quad.\quad.\quad Die Anordnung des Calculs der Störungen nach dem gewöhnlichen Verfahren hatte ich 1802, als ich es auf die Ceres anwandte, mir selbst entwickelt; allein meine Papiere sind freilich nicht so, dass ein anderer als ich sie gebrauchen könnte. Indessen müssen die Vorschriften von \textsc{Laplace} in der Mec. Cel. Vol. I oder auch von \textsc{Schubert} B. 3 im Wesentlichen damit übereinstimmen.\quad.\quad.\quad.\quad.\quad.
\par\medskip
\begin{center}\rule{.27\linewidth}{.4pt}\end{center}
